\documentclass[11pt]{article}
\usepackage{graphicx} % This lets you include figures
\usepackage{hyperref} % This lets you make links to web locations
\usepackage[margin=0.5in]{geometry}
\usepackage[rightcaption]{sidecap}
\usepackage{subcaption}
\usepackage{wrapfig}
\usepackage{float}
\usepackage{imakeidx}
\usepackage{indentfirst}
\makeindex
%---------------------------Do Not Edit Anything Above This Line!!------------------------

% edit the line below, if needed, to change the directory name for your image files.
\graphicspath{ {./images/} }



\begin{document}

%---------------------------Edit Content in the Box to Create the Title Page--------------
\begin{titlepage}
   \begin{center}
       \vspace*{1cm}
	   \Huge
       \textbf{Smart Resume/Job Search App}

       \vspace{0.5cm}
       \Large
       Client: Dr. Galloway \\
       Sprint Number 2\\
       9/30/2025 \\
   \end{center}

       \vspace{1.5cm}

\begin{table}[h!]
\centering
\begin{tabular}{|l|l|}
\hline
\textbf{Name} & \textbf{Email Address} \\ \hline
Brendan Rucker & brendan.rucker964@topper.wku.edu \\ \hline
Yukai Yang  & yukai.yang457@topper.wku.edu         \\ \hline
Michael Goldsmith & michael.goldsmith370@topper.wku.edu     \\ \hline
\end{tabular}
\end{table}

%Latex Table Generator    
%https://www.tablesgenerator.com/     
        
\vspace{4in}

\centering  
CS 360 \\
Fall 2025\\
Project Technical Documentation

\end{titlepage}
%---------------------------Edit Content in the Box to Create the Title Page--------------


% No text here.


%---------------------------Do Not Edit Anything In This Box!!------------------------
%Table of contents and list of figures will be autogenerated by this section.
\newpage
\setcounter{page}{1}%
\cleardoublepage
\pagenumbering{gobble}
\tableofcontents
\cleardoublepage
\pagenumbering{arabic}
\clearpage
\newpage
\setcounter{page}{1}%
\cleardoublepage
\pagenumbering{gobble}
\listoffigures
\cleardoublepage
\pagenumbering{arabic}
\newpage
%---------------------------Do Not Edit Anything In This Box!!------------------------




%---------------------------Project Introduction Section------------------------------

% No text here.

\section{Introduction} %\section{} is used to create major section headers

% No text here.

%---------------------------Project Overview------------------------------------------
\subsection{Project Overview} %\subsection{} is used to create minor sections 
% 300 words
% Description of the project, what the project provides, its purpose, problems solved, and target audience.

The Smart Resume and Job Search Application is designed to provide job seekers with a powerful tool that leverages AI to improve their employment prospects. At its core, the system uses natural language processing and content optimization techniques to analyze resumes, identify key strengths, and highlight areas for improvement. By parsing resumes in common formats such as PDF, DOCX, and TXT, the application ensures accessibility for a wide range of users and simplifies the uploading process.
%use blank lines to begin a new paragraph

The main purpose of this project is to address the common problems job seekers face when preparing resumes and searching for positions. Many applicants struggle with tailoring their resumes for specific job postings, often missing crucial keywords or failing to align their skills with employer expectations. This system solves that problem by matching the extracted resume data with job listings from integrated APIs like LinkedIn or Indeed, or through locally stored job datasets for offline use. Once matches are found, the application generates tailored resumes optimized for each job posting, significantly increasing the likelihood of catching the attention of recruiters.
%use blank lines to begin a new paragraph

The target audience for this project includes recent graduates, professionals seeking career changes, and individuals reentering the workforce. Ultimately, the application empowers users to present themselves more effectively, streamlines the job search process, and bridges the gap between job seekers and employers.

%---------------------------End Project Overview---------------------------------------

% No text here.

%---------------------------Project Scope----------------------------------------------
\subsection{Project Scope}
% 350 words
% Description of all deliverables, benefits, outcomes, and work required (all tasks, costs, time, people, resources, dates/deadlines, and final deliverables date).

The Smart Resume and Job Search Application will be delivered as a complete desktop system that enables resume analysis, customization, and job matching through AI-driven natural language processing. The final deliverables include: (1) a desktop application for Windows, (2) an integrated resume parser supporting PDF, DOCX, and TXT formats, (3) an AI-powered resume optimization module, (4) a job matching system linked to public job board APIs and offline datasets, (5) export functionality for resumes and reports, (6) a secure database for user data and job listings, and (7) performance, security, and user interface metric dashboards.

The project provides benefits such as faster resume tailoring, increased chances of passing ATS filters, streamlined job searches, and data security through encryption. Expected outcomes include a fully functional tool capable of processing a two-page resume within five seconds, offering both online and offline functionality, and empowering job seekers to efficiently connect with employers.

Project costs are limited to development resources, software tools, and potential API access fees. The team consists of 3 developers. Resources required include development machines, test environments for Windows OS, and job board API access.

The timeline spans 12 weeks (four sprints). Sprint 1 concludes at week 3, Sprint 2 at week 6, Sprint 3 at week 9, and Sprint 4 at week 12. The final deliverable is due at the end of Sprint 4, with a completed, tested, and documented system ready for demonstration.

%---------------------------End Project Scope---------------------------------------

% No text here.


\subsection{Technical Requirements}


%---------------------------Functional Requirements----------------------------------------------
\subsubsection{Functional Requirements} %\subsubsection{} used to create sections for parent subsections.
% Functional requirements define what a system or software must do, specifying the desired behavior or functionality.

% List as atomic bullet points that can be tested

\begin{table}[h!]
\centering
\begin{tabular}{|p{15cm}|}
\hline
\textbf{Mandatory Functional Requirements} \\ \hline
System will use an AI model or framework to perform natural language processing and content optimization. \\ \hline

The system will allow users to upload resumes in common formats (e.g., PDF, DOCX, TXT) for parsing and analysis.\\ \hline

The system will analyze the resume content to identify skills, education, work experience, and keywords, and provide feedback.\\ \hline

The system will match the extracted resume data with job listings obtained from integrated public job board APIs or from locally stored datasets provided by the user.\\ \hline

The system will generate tailored resume versions optimized for specific job postings, incorporating relevant keywords and formatting adjustments.\\ \hline
The system will allow users to export feedback reports, tailored resumes, and job match results in downloadable formats\\ \hline

\end{tabular}
\end{table}

% Paragraph (150 words) explaining the need and purpose for the listed Functional Requirements.
The functional requirements of the Smart Resume and Job Search Application are essential for ensuring the system meets user needs and achieves its intended purpose. By integrating an AI model capable of natural language processing and content optimization, the system addresses the common challenge of resumes not being tailored effectively for specific job postings. Supporting uploads in standard formats like PDF, DOCX, and TXT ensures accessibility and ease of use for all users. Parsing and analyzing resume content to extract skills, education, and experience provides the foundation for generating meaningful feedback and aligning resumes with employer expectations. The job matching feature, which connects extracted data with public job board APIs or offline datasets, directly benefits users by identifying opportunities that fit their profiles. Finally, enabling tailored resume generation and export in professional formats ensures that applicants are well-prepared, organized, and competitive in today’s technology-driven job market.

%---------------------------End Functional Requirements----------------------------------------------

% No text here.

%---------------------------Non-Functional Requirements----------------------------------------------
\subsubsection{Non-Functional Requirements}
% Non-functional requirements specify the constraints, qualities, or attributes that the system or software must possess, such as performance, security, usability, portability, fault tolerance, or reliability.

% List as atomic bullet points that can be tested

\begin{table}[h!]
\centering
\begin{tabular}{|p{15cm}|}
\hline
\textbf{Mandatory Non-Functional Requirements} \\ \hline
The system will process and return resume analysis results within 5 seconds for a standard two-page resume.\\ \hline

The system will operate on Windows\\ \hline

The system will store all user data and generated outputs securely using encryption during storage and transmission.\\ \hline

The system will maintain compatibility with the latest major version of each integrated job board API.\\ \hline

The system will allow for offline operation when job matching is performed using locally stored job listing datasets.\\ \hline

All project source code must be developed by the CS 360 project team.\\ \hline

The project must use a database.\\ \hline

Performance metrics should be gathered and optimized.\\ \hline

Security metrics should be gathered and optimized.\\ \hline

User interface metrics should be gathered and optimized.\\ \hline

\end{tabular}
\end{table}

% Paragraph (150 words) explaining the need and purpose for the listed Non-Functional Requirements.
The non-functional requirements of the Smart Resume and Job Search Application are critical to ensuring that the system not only works correctly but also performs reliably and securely. Fast processing within five seconds for a two-page resume ensures efficiency and keeps users engaged, addressing the expectation of quick feedback in modern applications. Secure storage and encrypted data transmission are necessary to protect sensitive personal information such as resumes and job history from unauthorized access. Maintaining compatibility with the latest job board APIs guarantees the system’s relevance and functionality over time. Offline operation provides flexibility, allowing users to continue job searches without constant internet access. Finally, gathering performance, security, and user interface metrics ensures continuous improvement, optimization, and usability. Together, these requirements provide stability, trustworthiness, and long-term sustainability for the application.

%---------------------------End Non-Functional Requirements---------------------------------------

% No text here.

%---------------------------Target Hardware Details----------------------------------------------
\subsection{Target Hardware Details}
% 200 words
% CPU (if a specific architecture is needed), RAM (required while the product is running), Persistent Storage Space, Network connection (Wi-Fi, Ethernet), Network bandwidth (required while the product is running), Output devices (Monitor (how many? What resolution?), speakers, VR headset), Input devices (Keyboard, mouse, touchscreen, VR headset).  Create test cases for each to verify.
The Smart Resume and Job Search Application is designed to operate on standard desktop and laptop systems on Windows OS. A multi-core CPU with x86-64 architecture is recommended to ensure fast resume parsing and AI-driven analysis. At least 8 GB of RAM is required while the product is running to handle resume processing, job matching, and UI interactions smoothly. Persistent storage of 500 MB is needed for installation, user data, logs, and temporary cache files. A stable network connection, either Wi-Fi or Ethernet, is necessary for accessing public job board APIs; offline operation requires only local datasets. Minimum network bandwidth of 5 Mbps ensures timely retrieval of job postings and secure synchronization. Output devices include at least one monitor . Input devices required include a keyboard and mouse. VR headsets are not required.

%---------------------------End Target Hardware Details----------------------------------------------

% No text here.

%---------------------------Software Product Deveopment----------------------------------------------
\subsection{Software Product Development}
% 200 words
% List the following used with their purpose for development: IDEs, IDE plugins, Software Languages, Software Frameworks, Version Control, Asset Generation Tools, and tools for colaboration (Google Drive, One Drive, GitHub, etc.).  Describle how each tool is used within your team.
For development of the Smart Resume and Job Search Application, the team will utilize a variety of tools to streamline coding, collaboration, and project management. The primary IDE is Visual Studio Code. 

The core programming languages are Python for natural language processing and AI integration, and JavaScript for backend and UI components.

For version control, the team uses GitHub, enabling source code tracking, branching, merging, and collaborative development.

For collaboration, Google Drive is used to share and co-edit project documents, reports, and presentations. Discord provides real-time communication for daily discussions, sprint planning, and quick troubleshooting. Together, these tools create an efficient ecosystem for development, testing, and delivery of the project.

%---------------------------End Software Product Deveopment-----------------------------------------

% No text here.

%---------------------------End Project Introduction Section----------------------------------------


% No text here.





%---------------------------Modeling and Design Section------------------------------
\section{Modeling and Design}
% No text here.


%---------------------------System Boundaries------------------------------
\subsection{System Boundaries}

\subsubsection{Physical}
%100 words
% Describe the Physical System Boundaries.
For our Physical System Boundaries, we need to have the application running on a device that supports Windows. The device will also need to be able to have a network connect, either from ethernet or WiFi. Most of the components will be ran locally, but a network connection is needed to communicate to Job Search APIs. A local database server would have to be established on this device as well in order for local data storage as well as caching job searches for local use in offline mode. 

%uncomment the section below when you're ready to insert an image
%\begin{figure}[h!]
%    \centering
%    \includegraphics[width=1\textwidth]{images/picture_of_Physical_boundaries.png}
%    \caption{Description of the image here.}
%    \label{phy_boundaries}
%\end{figure}






\subsubsection{Logical}
%100 words
% Describe the Logical System Boundaries.
For our Logical System Boundaries, we need to ensure each component is able to clearly communicate with the component it needs to communicate with. Most of the components in our system will only need to communicate with one or two other components since we will have low cohesion in our system. Components are generally ran sequentially and only need to communicate with the prior and future step. We will need Job Searching APIs that will find job data that we store in our local database. For our components, we need a Resume Parser, a Job Searcher, a Resume Optimizer, and components that will Store and Recall data from the database. 

%uncomment the section below when you're ready to insert an image
%\begin{figure}[h!]
%    \centering
%    \includegraphics[width=1\textwidth]{images/picture_of_Logical_boundaries.png}
%    \caption{Description of the image here.}
%    \label{log_boundaries}
%\end{figure}

%---------------------------End System Boundaries------------------------------

% No text here.

%---------------------------Wireframes and Storyboard----------------------------------
\subsection{Wireframes and Storyboard}
%150 words
% Use wireframes to create scenes and images of the user interface.  Connect the wireframes to show progression through the product to create a storyboard.  Describe the wireframes and storyboard.
Text goes here.

%uncomment the section below when you're ready to insert an image
%\begin{figure}[h!]
%    \centering
%    \includegraphics[width=1\textwidth]{images/picture_of_storyboard.png}
%    \caption{Description of the image here.}
%    \label{storyboard}
%\end{figure}

%---------------------------End Wireframes and Storyboard----------------------------------

% No text here.

%---------------------------Unified Modeling Language----------------------------------
\subsection{UML}

\subsubsection{Class Diagrams}
%At least one for each design pattern category.   
%Each class diagram should include the following:
%	Title for each Class Diagram
%	Description of how and why each class is used
%	Mapping to source code (in the Appendix, specific files and line numbers) of where it’s used.

%uncomment the section below when you're ready to insert an image
\begin{figure}[h!]
    \centering
    \includegraphics[width=1\textwidth]{Images/class-creational.png}
    \caption{Class Diagram w/ Creational Design Pattern}
    \label{class_diagram}
\end{figure}

For the Creational Class diagram, there are 5 core components: ResumeParser, JobSearch, Optimizer, StoreData, and RecallData. To separate communication between these components and the client, a CoreFactory is established to create each of these components. The ResumeParser has it's own factory in order to create new ResumeParsers based on the file type of the inputted Resume. Currently, we only expect to support PDF, DOCX, and TXT file types. JobSearch also has a factory in order to create different API communication elements, one for each API. We currently have only found Google's Job Search API as a usable API for this task. We have a few singletons as well, starting with the AIModel singleton. This is to ensure the AI model is instantiated upon booting the application and isn't needed to be recreated later as the AI communication components will only ever need to communicate with the same instance; this is to reduce our use of resources to ensure faster responses. The Database also has a singleton for the same reason; we only need 1 instance of the database as it needs to save all of the user's data, even upon shutdown. The StoreData and RecallData components are connected to this Database Singleton for input and output communication with the Database. 

\begin{figure}[h!]
    \centering
    \includegraphics[width=.5\textwidth]{Images/class-behavioral.png}
    \caption{Class Diagram w/ Behavioral Design Pattern}
    \label{class_diagram}
\end{figure}

The Behavioral Class Diagram is focused on watching for state changes in our system. We will do this by using the ResumeParser to communicate with an Observer, what will then update the state based on a JobMatcher, ResumeCustomizer, and UI components. 

\begin{figure}[h!]
    \centering
    \includegraphics[width=.3\textwidth]{Images/class-structural.png}
    \caption{Class Diagram w/ Structural Design Pattern}
    \label{class_diagram}
\end{figure}

The Structural Class Diagram has the client communicating with a ResumeAppFacade class, which will be used to remove communication from the client and instead direct it from the Facade instead. This Facade is creates two components, a ResumeManager and a JobRecommender. The ResumeManager will parse and store resume data while the JobRecommender will get data by creating a JobSource, which uses the JobAPIAdapter. The JobAPIAdapter creates an ExternalAPI, which calls the Job APIs, which will overall grab job data and return it back to the JobRecommender, which can then be used for future classes. 

\subsubsection{Use Case Diagrams}
%Enough to cover all technical functional requirements.
%Each Use Case Diagram should include the following:
%	Title for each Use Case Diagram
%	A description of information given in the diagram.
%	Mapping to source code (in the Appendix, specific files and line numbers) of where it’s used.
The Use Case Diagram shows that our user will have a total of 6 use cases for this application. They will Upload their Resume, Analyze Resume Content, Get Feedback and Suggestions, View Job Recommendations, Generate Tailored Resumes, and Export Resume/Reports. The View Job Recommendations use case will also communicate with a Job Board API and a Local Dataset (if user is offline). 

%uncomment the section below when you're ready to insert an image
\begin{figure}[h!]
    \centering
    \includegraphics[width=.8\textwidth]{Images/use-case.png}
    \caption{Use Case Diagram}
    \label{use_case_diagram}
\end{figure}


\subsubsection{Use Case Scenarios Developed from Use Case Diagrams (Primary, Secondary)}
%Should include at least one primary and zero to many secondary scenarios
%Each Use Case Scenario should include the following:
%	Title for each Use Case Scenario
%	A short description of the information given in the scenario.
%	Mapping to source code (in the Appendix, specific files and line numbers) of where it’s used
Use Case 1: Upload Resume

- Actors: User

- Preconditions: User has a resume file in PDF, DOCX, or TXT format.

- Postconditions: The resume is uploaded and stored for further analysis.

- Main Flow:
User selects the "Upload Resume" option.
System prompts the user to choose a file from their device.
User selects a valid resume file.
System uploads the file and confirms the upload.

- Alternate Flow:
1a. If the file format is unsupported → system displays an error and asks the user to retry.
1b. If the file is corrupted → system notifies the user.

Use Case 2: Analyze Resume Content

- Actors: User, NLP Engine (system service)

- Preconditions: A resume has been uploaded successfully.

- Postconditions: Resume content is parsed and structured (skills, education, work experience, keywords).

- Main Flow:
User chooses "Analyze Resume."
System sends the resume to the NLP engine.
NLP engine extracts structured data (skills, education, experience, keywords).
System stores and displays the extracted data to the user.

- Alternate Flow:
2a. If parsing fails → system shows error message and allows retry.

Use Case 3: Get Feedback And Suggestions

- Actors: User, AI Optimization Engine

- Preconditions: Resume has been analyzed.

- Postconditions: User receives improvement suggestions.

- Main Flow:
User requests feedback.
System runs AI-based content optimization.
System generates suggestions (e.g., missing keywords, formatting improvements).
System displays a feedback report.

- Alternate Flow:
2a. If AI engine fails → system notifies user and suggests retry later.

Use Case 4: View Job Recommendations

- Actors: User, Job Board API, Local Dataset

- Preconditions: Resume has been uploaded and analyzed.

- Postconditions: User sees a ranked list of job postings relevant to their profile.

- Main Flow:
User selects "View Job Recommendations."
System queries external job board APIs.
If offline, system queries local dataset.
System matches resume keywords with job requirements.
System ranks and displays job recommendations.

- Alternate Flow:
2a. If APIs are unavailable → system switches to offline dataset.
2b. If no matches are found → system suggests general job categories.

Use Case 5: Generate Tailored Resume

- Actors: User, AI Optimization Engine

- Preconditions: Job recommendation has been selected.

- Postconditions: A customized resume is generated for that specific job.

- Main Flow:
User chooses a job from recommendations.
System creates a tailored resume using AI optimization.
System highlights inserted keywords and adjusted sections.
Tailored resume is saved for export.

- Alternate Flow:
2a. If tailoring fails → system falls back to the original resume with minor edits.

Use Case 6: Export Resume/Reports

- Actors: User, Export Service

- Preconditions: Resume analysis/feedback or tailored resume has been generated.

-Postconditions: User downloads a file in the chosen format (PDF or DOCX).

- Main Flow:
User clicks "Export Resume" or "Export Report."
System prompts for format selection (PDF/DOCX).
User selects the desired format.
System generates the file.
User downloads the exported file.

- Alternate Flow:
4a. If file generation fails → system displays error and allows retry.

%uncomment the section below when you're ready to insert an image
%\begin{figure}[h!]
%    \centering
%    \includegraphics[width=1\textwidth]{images/use_case_scenario.png}
%    \caption{Description of the image here.}
%    \label{use_case_scenario}
%\end{figure}


\subsubsection{Sequence Diagrams}
%Each diagram should include all actors/resources to cover one Use Case Scenario.
%	Each Sequence Diagram should include the following:
%	Title for each Sequence Diagram
%	A short description of the information given in the diagram.
%	Mapping to source code (in the Appendix, specific files and line numbers) of where it’s used
For our Sequence Diagram, we have our user (A Job Seeker) and 6 components: UI, Controller, Job Search APIs, Database, Parsing AI, and Generative AI. The UI is what displays content to the Job Seeker. The Controller is the main component that runs the application. The Database is a local database within the application. 

In our sequence, we start with the Job Seeker booting up the application. Here, we do a check to see if they are online; if they are we will do a quick location-only based job search through the Job Search APIs and stash the data as an initial source for the local database if the user goes offline. The Controller then waits for the Job Seeker to input their resume. Upon doing so, the Resume is sent to the Parsing AI where we grab all important information like skills, education, previous work experience, etc. This data is then sent back to the Controller and sent to the Job Searching APIs which find relevant jobs based on this dataset. Upon returning a list of jobs, it will output the options back to the Job Seeker to allow them to select which job they want to pursue. Upon a selection, the information and their Resume is sent to the Generative AI model that will then optimize the Resume's text based on the job description. All data is then saved for the Job Seeker to view later. At each step in this process, the Controller is also signaling state changes to the UI, which updates the UI's display. 

%uncomment the section below when you're ready to insert an image
\begin{figure}[h!]
    \centering
    \includegraphics[width=.7\textwidth]{Images/sequence.png}
    \caption{Sequence Diagram}
    \label{use_case_scenario}
\end{figure}

\subsubsection{State Diagrams}
%Each diagram should identify the object/resource/asset (in the title of the diagram) and display all states and actions/events that create state changes.
%	Each State Diagram should include the following:
%	Title for each State Diagram
%	A short description of the information given in the diagram
%	Mapping to source code (in the Appendix, specific files and line numbers) of where it’s used
In this State Diagram, we can see there are many states that change throughout the process of this application, including: Idle, Uploading, Parsing, Searching, Matching, Optimizing, and Exporting. In the Idle state, this is when we are waiting for an input from the user. The Uploading state is the process in which the resume is currently being received and ends when it is sent to the Parsing state, which grabs data from the Resume. Searching is when data from the Parsing State is sent to the Job Searching APIs. When that data is received, it moved to the Matching State which matches the skills from the resume to the job's provided. We then move onto the Optimizing State when the user selects the job they want to pursue and begins the analysis to specialize their resume to that particular job. The Exporting State is when optimization is complete and the user can choose which file type to export the document in. 

%uncomment the section below when you're ready to insert an image
\begin{figure}[h!]
    \centering
    \includegraphics[width=.8\textwidth]{Images/state.png}
    \caption{State Diagram}
    \label{state_diagram}
\end{figure}

\subsubsection{Component Diagrams}
%Single diagram that defines component APIs and communication pathways for all software and dependencies used in the product.
%	The Component Diagram should include the following:
%	Title for the Component Diagram
%	A description of the information given in the diagram
%	Protocols used for each communication pathway
%	Mapping to all source code and dependency software (in the Appendix, specific files)
For the Component Diagram, we have a User Interface, a Resume Parser, an AI/NLP Analysis, a Job Matching, Job Search APIs, a Resume Customization, and an Export component. The User Interface communicates to the Resume Parser when it has received a Resume input, when then communicates to the AI Analysis component to gather data from that resume. That data is then communicated to the Job Search APIs, which then sends it to the Job Matching component. That data is returned to the User Interface, which communicates to the Resume Customization Component when a job is selected. The process ends with communication to the Export Component. 

%uncomment the section below when you're ready to insert an image
\begin{figure}[h!]
    \centering
    \includegraphics[width=1\textwidth]{Images/component.png}
    \caption{Component Diagram}
    \label{component_diagram}
\end{figure}

\subsubsection{Deployment Diagrams}
%Single diagram that defines the physical nodes, virtual nodes, and communication for all software related to the product.  Should be an extension of the Component Diagram.
%	The Deployment Diagram should include the following:
%	Title for the Deployment Diagram
%	A description of the information given in the diagram
%	Identify all physical and virtual nodes used for process execution
%	Identify all internal and external node communication
In our Deployment Diagram, we have a Physical Application Server and a Database Server. All of these are within the same physical device that is the client. All of the components within the component diagram are within the application server, with no outgoing connections besides the Job Search APIs component, which does sent outgoing communication to those APIs somewhere within the internet. 

%uncomment the section below when you're ready to insert an image
\begin{figure}[h!]
    \centering
    \includegraphics[width=.6\textwidth]{Images/deployment.png}
    \caption{Deployment Diagram}
    \label{deployment_diagram}
\end{figure}


%---------------------------End Unified Modeling Language----------------------------------

%---------------------------Requirments Traceability------------------------------------------------
\subsection{Requirements Traceability Table}
%150 words and table
%	Create a table that compares all requirements to use cases to make sure there are no missing features.
\begin{table}[]
    \centering
    \begin{tabular}{|c|c|c|c|c|c|c|}
        & Upload Resume & Analyze Resume & Get Feedback & View Jobs & Optimize Resume & Export Resume \\
        Natural Language Processing & X & X & X & & X & \\
        Job Searching & X & & & X & & X \\
    \end{tabular}
    \caption{Requirements Traceability Table}
    \label{tab:placeholder}
\end{table}
%------------------------------------------------------------------------------------

% No text here.

%---------------------------Version Control------------------------------------------------
\subsection{Version Control}
%150 words
%Describe your group's approach to version control.  Include details on how and when you make commits to GitHub.  Explain version naming, branching, and integration approaches.
Text goes here.


%---------------------------End Version Control------------------------------------------------

% No text here.




%---------------------------End Data Dictionary------------------------------------------------


% No text here.

%---------------------------Data Dictionary------------------------------------------------
\subsection{Data Dictionary}
%150 words and table
%	Create a table that displays your Data Dictionary and describe how it is being used to define data structures and other major variables/elements in the software product.
Text and table goes here.


%---------------------------End Data Dictionary------------------------------------------------


% No text here.


%---------------------------User Experience Details------------------------------------------------
\subsection{User Experience}

\subsubsection{Gameplay Diagram}
%100 words
%Also called a controlled flow graph.  Insert an image of the gameplay diagram and describe the diagram in detail.
Text goes here.

%uncomment the section below when you're ready to insert an image
%\begin{figure}[h!]
%    \centering
%    \includegraphics[width=1\textwidth]{images/gameplay_diagram.png}
%    \caption{Description of the image here.}
%    \label{gameplay_diagram}
%\end{figure}

\subsubsection{Gameplay Objectives}
%100 words
%Should explain why someone should play your game. Should define the goals of your game.
Text goes here.

\subsubsection{User Skillset}
%100 words
%Describe the abilities a user should have in order to interact with the software: Screen tapping, Memorization, Puzzle, solving, Strategy development, Quick reflexes 
Text goes here.

\subsubsection{Gameplay mechanics}
%100 words
%Thorough guide of the game rules: Turn-based, Random chances, Capturing/eliminating, Time-based, Team-based
Text goes here.

\subsubsection{Gameplay Items}
%100 words
%This section of the GDD refers to smaller elements of the game, such as: Health/power-ups, Coins, Loot crates, Weapons/character upgrades, Etc. Changes gameplay elements as the player progresses through the game.
Text goes here.

\subsubsection{Gameplay Challenges}
%100 words
%Description of game challenges as a player progresses. Identify how a player handles difficulties and ensures the player can obtain essential tools to continue progressing through the game.
Text goes here.

\subsubsection{Gameplay Menu Screens}
%100 words and wireframes
%Game start menu screen, pause menu screen, game over screen
%Give wireframe/drawing/image for all menus in the game.
Text goes here.

\subsubsection{Gameplay Heads-Up Display}
%100 words and wireframes
%Current Game information visually relayed to the player as a part of the game’s user interface.
%HUD features are generally static, on-screen so they remain visible during gameplay. Health/Lives, Time, Weapons/Ammo, Capabilities, Menus, Mini-Map, Score, etc.
Text goes here.

\subsubsection{Gameplay Art Style}
%50 words
%describe the visual appearance of characters, objects, and other in-game elements. (Realistic, Cartoon, etc.)
Text goes here.


\subsubsection{Gameplay Audio}
%100 words
%Description of all game audio: Character sound/voices, Ambient world sounds, Background music, etc.
Text goes here.



%---------------------------End User Experience Details------------------------------------------------

% No text here.


%---------------------------End Modeling and Design Section----------------------------------


% No text here.


%---------------------------Non-Functional Product Details Section---------------------------
\section{Non-Functional Product Details}

%---------------------------Product Security-------------------------------------------------

\subsection{Product Security}

\subsubsection{Approach to Security in all Process Steps}
%200 words
%Describe how your team modified the original technical document to address security issues in the Requirements, Modeling and Design, and Implementation sections.
Text goes here.


\subsubsection{Security Threat Model}
%150 words and Security Threat Model
%Create and add a Security Threat Model, related to the Deployment Diagram, that identifies trust boundaries and potential security risks
Text goes here.

\subsubsection{Security Levels}
%150 words
%Describe the different security levels for general users and administrators.  Also describe the authentication/authorization techniques for users of the software product.
Text goes here

%---------------------------End Product Security-------------------------------------------------

% No text here.

%---------------------------Product Performance-------------------------------------------------

\subsection{Product Performance}


\subsubsection{Product Performance Requirements}
%100 words and list of performance non-functional requirements.
%Define and justify performance requirements.  These should be added to the list of non-functional requirements.
Text goes here.

\subsubsection{Measurable Performance Objectives}
%100 words and list of objectives.
%Measurable Performance Objectives should be stated in this section that relate to the performance requirements.
Text goes here.

\subsubsection{Application Workload}
%200 words
%Application Workload information should be gathered and visualized in this section.  This generally requires historical data on how the software product is being used.  For example, users generally spend 10% of the time interacting with menus, 80% of the time interacting with main features, 5% of the time saving work, etc.  These workloads should not be assumptions or guesses.  Timers need to be created for all major UI features of the product to generate reliable application workload analysis.
Text goes here.

\subsubsection{Hardware and Software Bottlenecks}
%200 words
%Hardware and Software Bottlenecks should be identified and discussed in this section, with test cases to justlify.
Text goes here.

\subsubsection{Synthetic Performance Benchmarks}
%250 words
%Synthetic Performance Benchmark test cases should be developed and executed on target hardware.  Results should be visualized and discussed in terms of the required target hardware details. (File I/O, CPU, Database).  Sysbench
Text goes here.


\subsubsection{Performance Tests}
%250 words and test case description with results
%Performance Test Cases should be given in this section.  Test Cases should be executed and results should be visualized (Images, Graphs, etc.)
%examples of performance tests include, but not limited to: Load testing (expected and peak loads, exceeding peak loads), Stress testing, throughput testing, function call timers, compatability testing, fault tolerance testing, etc.
Text goes here


%---------------------------End Product Performance-------------------------------------------------

% No text here.


%---------------------------End Non-Functional Product Details Section---------------------------



% No text here.



%---------------------------Software Product Testing Section-------------------------------------
\section{Software Testing}



%---------------------------Software Testing Plan Template-------------------------------------

\subsection{Software Testing Plan Template}
%Each of the testing levels (unit, Integration, System, Acceptance) should use the following test plan template.

\textbf{Test Plan Identifier:} %Provides a unique identifier for the test. Every test should have a unique identification number for reference.

\textbf{Introduction:} % 50 words. Brief description and objective about the test type.

\textbf{Test item:} %50 words. Includes detailed information about the Software Under Test (SUT).

\textbf{Features to test/not to test:} %50 words. In scope features. This could be newly added or updated features. Out of scope features not tested. [Provide reasoning for exclusion, like, non-impacted, low priority, etc.]

\textbf{Approach:} %50 words. Strategy to test the software. Includes types of tests and how to test. Functional, performance, security testing using combined [manual + automation], manual only, automation only approach.

\textbf{Test deliverables:} %50 words. All the deliverables from the testing e.g. approaches, test cases, reports etc.

\textbf{Item pass/fail criteria:} %50 words. Entry and Exit criteria for all items. 

\textbf{Environmental needs:} %50 words. Infrastructure required for SUT and executing test cases.

\textbf{Responsibilities:} %50 words. Roles and responsibilities for various testing / supported activities.

\textbf{Staffing and training needs:} %50 words. Training needs to bridge the gap of available and expected skill.

\textbf{Schedule:} %50 words.  Test schedule should also be noted in the Gantt Chart. Test estimation (Efforts) and high-level schedule. Schedule should be for key deliverables or important milestones. Ideally, all test deliverables included in the test plan should be scheduled.

\textbf{Risks and Mitigation:} %100 words. Risk identification for applicable items, assumptions, and mitigation plan.

\textbf{Approvals:} %Approvals and sign of dates.

%---------------------------Software Testing Plan Template-------------------------------------


% No text here.



%---------------------------Unit Testing-------------------------------------
\subsection{Unit Testing}
%copy of the completed Unit test plan should be placed here.
Text goes here.

\subsubsection{Source Code Coverage Tests}
%100 words
%Insert Flow Graph Image(s).  Define cyclomatic complexity, basis paths, and Unit Test Cases
Text goes here.

%uncomment the section below when you're ready to insert an image
%\begin{figure}[h!]
%    \centering
%    \includegraphics[width=1\textwidth]{images/flow_graph.png}
%    \caption{Description of the image here.}
%    \label{flow_graph}
%\end{figure}



\subsubsection{Unit Tests and Results}
%All unit tests results visualized (table, graph, etc.)
Text goes here.


%---------------------------End Unit Testing-------------------------------------

% No text here.

%---------------------------Integration Testing-------------------------------------
\subsection{Integration Testing}
%copy of the completed Integration test plan should be placed here.
Text goes here.

\subsubsection{Integration Tests and Results}
%All integration tests results visualized (table, graph, etc.)
Text goes here.


%---------------------------End Integration Testing-------------------------------------

% No text here.


%---------------------------System Testing-------------------------------------
\subsection{System Testing}
%copy of the completed System test plan should be placed here.
Text goes here.

\subsubsection{System Tests and Results}
%All system tests results visualized (table, graph, etc.)
Text goes here.


%---------------------------End System Testing-------------------------------------

% No text here.


%---------------------------Acceptance Testing-------------------------------------
\subsection{Acceptance Testing}
%copy of the completed Acceptance test plan should be placed here.
Text goes here.

\subsubsection{Acceptance Tests and Results}
%All Acceptance tests results visualized (table, graph, etc.)
Text goes here.


%---------------------------End System Testing-------------------------------------

% No text here.





%---------------------------End Software Product Testing Section-------------------------------------


% No text here.



%---------------------------Conclusion Section-------------------------------------
\section{Conclusion}
%200 words
%Concluding remarks that summarizes the purpose and outcomes of the technical document.  Discussion of short comings and future work.
Text goes here.

%---------------------------End Conclusion Section-------------------------------------


% No text here.



%---------------------------Appendix Section-------------------------------------------
\section{Appendix}

\subsection{Software Product Build Instructions}
%Include in this section all steps for copying the current state of the product to new computers for continued development.
Text goes here.

\subsection{Software Product User Guide}
%Include in this section an overview guide on how to use your software product for a general user and an administrative user.
Text goes here.

\subsection{Source Code with Comments}
%Include in this section all final source code for the product.  Label each file with headings such as, C.1 file1.c, C.2 file2.c, C.3 file1.py, etc.  All source code should be effectively commented.
Text goes here.







%---------------------------End Appendix Section-------------------------------------------














%example image:  uncomment to show usage
%\begin{figure}[h]
%    \centering
%    \includegraphics[width=1\textwidth]{images/Add_non-music.png}
%    \caption{This is how you add non-music items.}
%    \label{fig16}
%\end{figure}


%example links:  uncomment to show usage.
%\url{https://www.youtube.com}
%\href{https://www.wku.edu/}{WKU Homepage}
%\footnote{You can put the link in a footnote like this.}

% Anything to the right of a percent sign will be ignored by LaTeX.
% You can use this to put notes to yourself.  



\end{document}
